\subsection{Matched Pairs and McNemar's Test}
\label{subsec:matched-pairs-and-mcnemars-test}
% Source integration: attached Chapter 13 screenshots, Section 13.5 and selected problems.

Matched categorical data require a different table from unmatched samples.  The pair, not the individual measurement, is the independent sampling unit, just as in the continuous paired-sample setting of Subsection \ref{subsec:paired-sample-inference}.

\begin{definition} \label{def:matched-pair-binary-table} \hlt{Matched Binary Pair Table}\\
For paired binary responses, write the table as
\[
\begin{array}{c|cc}
&\text{Second: }0&\text{Second: }1\\ \hline
\text{First: }0&n_{11}&n_{12}\\
\text{First: }1&n_{21}&n_{22}
\end{array}
\]
after choosing which level is coded as $0$ or $1$.  The diagonal cells count concordant pairs.  The off-diagonal cells count discordant pairs and contain the information about marginal change between the first and second members of each pair.
\end{definition}

\begin{definition} \label{def:mcnemars-null} \hlt{Marginal Homogeneity for Matched Pairs}\\
Let $\pi_{ij}$ be the cell probabilities in Definition \ref{def:matched-pair-binary-table}.  Equality of the two marginal response probabilities is
\[
\pi_{1\cdot}=\pi_{\cdot1}
\quad\text{and}\quad
\pi_{2\cdot}=\pi_{\cdot2},
\]
which is equivalent in a $2\times2$ matched table to
\[
H_0:\pi_{12}=\pi_{21}.
\]
Thus only the two off-diagonal probabilities are constrained by the null hypothesis.
\end{definition}

\begin{method} \label{meth:mcnemars-test} \hlt{McNemar's Test}\\
For the matched binary table of Definition \ref{def:matched-pair-binary-table}, the maximum likelihood estimates under $H_0:\pi_{12}=\pi_{21}$ are
\[
\widehat{\pi}_{11}=\frac{n_{11}}{n},
\qquad
\widehat{\pi}_{22}=\frac{n_{22}}{n},
\qquad
\widehat{\pi}_{12}=\widehat{\pi}_{21}=\frac{n_{12}+n_{21}}{2n}.
\]
The Pearson statistic reduces to the off-diagonal contrast
\[
X^2
=
\frac{(n_{12}-n_{21})^2}{n_{12}+n_{21}}.
\]
For large discordant-pair count $n_{12}+n_{21}$, this is calibrated by $\chi_1^2$.
\end{method}

\begin{example} \label{ex:hodgkins-tonsillectomy-matched-pairs} \hlt{Hodgkin's Disease and Sibling Controls}\\
The source contrasts an unmatched analysis with a matched sibling design.  When $85$ Hodgkin's patients were paired with disease-free siblings, the pair table was
\[
\begin{array}{c|cc}
&\multicolumn{2}{c}{\text{Sibling}}\\
\text{Patient}&\text{No tonsillectomy}&\text{Tonsillectomy}\\ \hline
\text{No tonsillectomy}&37&7\\
\text{Tonsillectomy}&15&26
\end{array}
\]
The ordinary unmatched chi-square comparison ignores the pairing.  McNemar's statistic uses only the discordant cells:
\[
X^2=\frac{(7-15)^2}{7+15}=2.91.
\]
The source reports the corresponding $p$-value as about $0.09$, which weakens the apparent evidence from the unpaired analysis.
\end{example}

\begin{example} \label{ex:cell-phone-driving-mcnemar} \hlt{Cell Phones and Driving}\\
In a matched case-crossover study, each driver involved in a collision was compared with the same driver's phone status at the same time during the previous week:
\[
\begin{array}{c|cc|c}
\text{Day of collision}&\multicolumn{2}{c|}{\text{Before collision}}&\\
&\text{On phone}&\text{Not on phone}&\text{Total}\\ \hline
\text{On phone}&13&157&170\\
\text{Not on phone}&24&505&529\\ \hline
\text{Total}&37&662&699
\end{array}
\]
McNemar's statistic is
\[
X^2=\frac{(157-24)^2}{157+24}=97.7,
\]
so the association is statistically strong.  The source also emphasizes the design limitation: the association need not mean that phone use directly causes collisions, because time-varying factors such as emotional stress could affect both phone use and driving risk.
\end{example}

\begin{remark} \label{rem:mcnemar-versus-independence} \hlt{Matched Does Not Mean Independent}\\
The null in McNemar's test is not independence of the two paired responses.  It is equality of marginal response probabilities within the pair.  Source problems ask for comparisons with paired $t$ procedures and with square matched tables; the common lesson is that the analysis must preserve the matching structure rather than treating the two margins as independent samples.
\end{remark}
